% Document/LinearAlgebra/VectorSpaces/SubspaceRepresentations/EquivalenceSimilarity.tex
% SEPARATE CHAPTER (draft) -- split out of the old combined Part 5. Equivalence (A = Q B P^{-1}) and
% similarity (A = Q^{-1} B Q) of maps/matrices; one-sided (domain/codomain) equivalence; matrices of one
% map are equivalent (similar for endomorphisms); kernel/range one-sided invariance; rank a complete
% equivalence invariant (forward ref); basis-free content = the class; det/trace/char.poly as forward
% refs. Depends on the Matrices chapter's invertibility section. Gaussian elimination / normal forms join
% this track later. DRAFT for review.

\newpage
\section{Equivalence and Similarity of Linear Maps}
\label{sec:equivalence-similarity}

\begin{intuition}
    The matrix of a fixed linear map $T : U \to V$ is not unique: it depends on the ordered bases chosen
    for source and target. So which matrices represent one and the same map? The change-of-bases formula
    (Corollary~\ref{cor:change-of-bases}) already answers in form --- two matrices of $T$ differ by
    invertible factors fore and aft --- and isolating that relation, abstracted from any particular $T$,
    produces the notion of \textbf{equivalence} of linear maps (hence of matrices). Tightening it to a
    \emph{single} change of basis on a self-map gives \textbf{similarity} of endomorphisms. The payoff is
    conceptual: the genuinely basis-free content of a map is precisely its equivalence class --- for an
    endomorphism, its similarity class --- so a quantity read off a matrix is an intrinsic invariant of
    the underlying map exactly when it depends only on that class. This chapter develops the two relations
    and their one-sided refinements; the \emph{classification} that completes the story --- normal forms,
    and the Gaussian elimination that computes them --- follows in the chapters to come.
\end{intuition}

Equivalence and similarity are, at bottom, relations between \emph{linear maps}: they say that two maps
agree after a change of coordinates --- on source and target independently for equivalence, by a single
change for similarity. We define them at the level of maps, where they are most transparent. There is no
separate notion for matrices: a matrix \emph{is} a linear map between coordinate spaces
(Section~\ref{sec:matrices-as-maps}), and the invertible matrices are exactly the automorphisms
$\GL{K^{(X)}}$ (Remark~\ref{rem:invertible-square}), so each definition below applies to matrices with
nothing added --- equivalence of matrices simply \emph{is} equivalence of the maps they are. We record
the resulting matrix form inline as each definition is given.

\begin{definition}[Equivalence of linear maps]
\label{def:map-equivalence}
    Linear maps $S, T : U \to V$ are \textbf{equivalent}, written $S \Equivalent T$, if there are
    isomorphisms $P \in \GL{U}$ and $Q \in \GL{V}$ with
    \[
        T = \Composition{Q}{\Composition{S}{P^{-1}}} ,
    \]
    that is, the square
    \begin{center}
    \begin{tikzcd}[column sep=large, row sep=large]
        U \arrow[r, "S"] \arrow[d, "P"'] & V \arrow[d, "Q"] \\
        U \arrow[r, "T"'] & V
    \end{tikzcd}
    \end{center}
    commutes ($\Composition{Q}{S} = \Composition{T}{P}$). One changes coordinates by $P$ in the source
    and, independently, by $Q$ in the target.

    \emph{In coordinates.} Taking $U = K^{(X)}$ and $V = K^{(Y)}$ --- so that $S, T$ are matrices
    $A, B \in \CFMat{Y}{X}[K]$ --- the very same definition reads: $A \Equivalent B$ iff there are
    invertible matrices $P \in \CFMat{X}{X}[K]$ and $Q \in \CFMat{Y}{Y}[K]$ with $A = Q B P^{-1}$, where
    $\circ$ is the matrix product. This \emph{is} the equivalence of matrices; it is not a fresh
    definition but the one above with coordinate spaces in place of $U, V$.
\end{definition}

Before tying the two ends together by a single change (similarity, below), it pays to isolate each end
on its own. Full equivalence reparametrises source and target at once; relaxing it to move only one at a
time splits $\Equivalent$ into two \emph{one-sided} relations, of which it is the composite.

\begin{definition}[Domain and codomain equivalence]
\label{def:one-sided-equivalence}
    Let $S, T : U \to V$ be linear maps. They are:
    \begin{itemize}
        \item \textbf{domain-equivalent} (or \textbf{column-equivalent}), written
              $S \DomainEquivalent T$, if $T = \Composition{S}{P^{-1}}$ for some $P \in \GL{U}$ --- only
              the source is reparametrised (the case $Q = \Identity{V}$ of
              Definition~\ref{def:map-equivalence});
        \item \textbf{codomain-equivalent} (or \textbf{row-equivalent}), written
              $S \CodomainEquivalent T$, if $T = \Composition{Q}{S}$ for some $Q \in \GL{V}$ --- only
              the target is reparametrised (the case $P = \Identity{U}$).
    \end{itemize}
    Each is again an equivalence relation, by Proposition~\ref{prop:equiv-is-equiv}.

    \emph{In coordinates.} For matrices $A, B \in \CFMat{Y}{X}[K]$ the same two definitions read:
    $A \DomainEquivalent B$ iff $A = B P^{-1}$ for an invertible $P \in \CFMat{X}{X}[K]$ --- right
    multiplication, remixing the \emph{columns} (indexed by the source $X$) --- and
    $A \CodomainEquivalent B$ iff $A = Q B$ for an invertible $Q \in \CFMat{Y}{Y}[K]$ --- left
    multiplication, remixing the \emph{rows} (indexed by the target $Y$). Again these are the map
    relations applied to coordinate spaces, not new notions.
\end{definition}

\begin{proposition}[Equivalence is a domain step followed by a codomain step]
\label{prop:equivalence-factors}
    For linear maps $S, T : U \to V$ the following are equivalent.
    \begin{enumerate}
        \item $S \Equivalent T$.
        \item There is a map $R : U \to V$ with $S \DomainEquivalent R$ and $R \CodomainEquivalent T$
              --- reparametrise the source, then the target.
        \item There is a map $R' : U \to V$ with $S \CodomainEquivalent R'$ and $R' \DomainEquivalent T$
              --- reparametrise the target, then the source.
    \end{enumerate}
    Hence, as composites of relations,
    \[
        \Equivalent \;=\; \CodomainEquivalent \circ \DomainEquivalent
                    \;=\; \DomainEquivalent \circ \CodomainEquivalent :
    \]
    the two one-sided changes commute, and either order realises full equivalence. The same holds
    verbatim for matrices, with $P, Q$ the invertible factors of $A = Q B P^{-1}$.
\end{proposition}

\begin{proof}
    $S \Equivalent T$ means $T = \Composition{Q}{\Composition{S}{P^{-1}}}$ with $P \in \GL{U}$,
    $Q \in \GL{V}$ (Definition~\ref{def:map-equivalence}). Associating the composite one way, set
    $R \defeq \Composition{S}{P^{-1}}$: then $S \DomainEquivalent R$ via $P$, and
    $T = \Composition{Q}{R}$ gives $R \CodomainEquivalent T$ via $Q$ --- this is~(2). Associating the
    other way, set $R' \defeq \Composition{Q}{S}$: then $S \CodomainEquivalent R'$ via $Q$, and
    $T = \Composition{R'}{P^{-1}}$ gives $R' \DomainEquivalent T$ via $P$ --- this is~(3). Conversely,
    substituting the intermediate map back into either chain reconstitutes
    $T = \Composition{Q}{\Composition{S}{P^{-1}}}$, so each of~(2),~(3) implies $S \Equivalent T$. The two
    factorisations agree because composition is associative,
    $\Composition{Q}{(\Composition{S}{P^{-1}})} = \Composition{(\Composition{Q}{S})}{P^{-1}}$, which is
    exactly the commutativity of the two one-sided relations.
\end{proof}

\begin{remark}[Reading the decomposition off the four-level diagram]
\label{rem:one-sided-diagram}
    The two one-sided equivalences are the two independent \emph{columns} of the four-level
    change-of-basis diagram (Remark~\ref{rem:active-passive-cob}): the source space $U$ with its frames
    occupies the left column, the target $V$ the right. Passing between the coordinate rows
    $\CoordMatrix[\mathcal B_U][\mathcal B_V]{T}$ and $\CoordMatrix[\mathcal C_U][\mathcal C_V]{T}$ is
    governed, on the right, by the source transition
    $\CoordMatrix[\mathcal C_U][\mathcal B_U]{\Identity{U}}$ and, on the left, by the target transition
    $\CoordMatrix[\mathcal B_V][\mathcal C_V]{\Identity{V}}$. Sliding down the \emph{source} column alone
    --- holding $\mathcal B_V = \mathcal C_V$ --- turns $\CoordMatrix[\mathcal B_U][\mathcal B_V]{T}$ into
    \[
        \CoordMatrix[\mathcal C_U][\mathcal B_V]{T}
        = \Composition{\CoordMatrix[\mathcal B_U][\mathcal B_V]{T}}
                      {\CoordMatrix[\mathcal C_U][\mathcal B_U]{\Identity{U}}} ,
    \]
    a domain (column) equivalence; sliding down the \emph{target} column alone gives
    $\CoordMatrix[\mathcal B_U][\mathcal C_V]{T}
    = \Composition{\CoordMatrix[\mathcal B_V][\mathcal C_V]{\Identity{V}}}{\CoordMatrix[\mathcal B_U][\mathcal B_V]{T}}$,
    a codomain (row) equivalence; performing both --- in either order --- is the full change-of-bases
    formula (Corollary~\ref{cor:change-of-bases}), hence the full equivalence
    $\CoordMatrix[\mathcal B_U][\mathcal B_V]{T} \Equivalent \CoordMatrix[\mathcal C_U][\mathcal C_V]{T}$
    (of the matrices, as the maps they are). The commutativity of
    Proposition~\ref{prop:equivalence-factors} is precisely the independence of the two columns: the
    source frame may be changed before or after the target frame, with the same result.
\end{remark}

\begin{definition}[Similarity of endomorphisms]
\label{def:map-similarity}
    Endomorphisms $S, T \in \End[K]{V}$ are \textbf{similar} (or \textbf{conjugate}), written
    $S \Similar T$, if there is an automorphism $Q \in \GL{V}$ with
    \[
        T = \Composition{Q^{-1}}{\Composition{S}{Q}} .
    \]
    Similarity is the special case of equivalence in which source and target are changed by the
    \emph{same} isomorphism ($U = V$, $P = Q$); it is precisely the conjugation action of $\GL{V}$ on
    $\End[K]{V}$.

    \emph{In coordinates.} For square matrices $A, B \in \CFMat{X}{X}[K]$ the same definition reads:
    $A \Similar B$ iff $A = Q^{-1} B Q$ for an invertible $Q \in \CFMat{X}{X}[K]$ --- once more the map
    relation in coordinate spaces, the conjugation orbit of $A$ under $\GL{K^{(X)}}$.
\end{definition}

\begin{proposition}[All four relations are equivalence relations]
\label{prop:equiv-is-equiv}
    Equivalence, column (domain) equivalence, and row (codomain) equivalence are equivalence relations on
    $\Hom[K]{U}{V}$, and similarity is an equivalence relation on $\End[K]{V}$.
\end{proposition}

\begin{proof}
    We prove it for $\Equivalent$ and obtain the other three as special cases. Reflexivity:
    $T = \Composition{\Identity{V}}{\Composition{T}{\Identity{U}^{-1}}}$, witnessed by the identity pair
    $(P, Q) = (\Identity{U}, \Identity{V})$. Symmetry: from $T = \Composition{Q}{\Composition{S}{P^{-1}}}$
    we get
    $S = \Composition{Q^{-1}}{\Composition{T}{P}} = \Composition{Q^{-1}}{\Composition{T}{(P^{-1})^{-1}}}$,
    witnessed by the inverse pair $(P^{-1}, Q^{-1})$. Transitivity: if
    $T = \Composition{Q}{\Composition{S}{P^{-1}}}$ and $S = \Composition{Q'}{\Composition{R}{(P')^{-1}}}$,
    then $T = \Composition{(\Composition{Q}{Q'})}{\Composition{R}{(\Composition{P}{P'})^{-1}}}$, witnessed
    by the composite pair $(\Composition{P}{P'}, \Composition{Q}{Q'})$, using
    $(\Composition{P}{P'})^{-1} = \Composition{(P')^{-1}}{P^{-1}}$ and that composites of isomorphisms are
    isomorphisms.

    The only facts used are that the automorphism pairs $(P, Q)$ contain the identity pair and are closed
    under componentwise inverse and composition. Each of the other three relations is $\Equivalent$ with
    $(P, Q)$ confined to a subset of $\GL{U} \times \GL{V}$ that is again closed under these three
    operations:
    \begin{itemize}
        \item $\CodomainEquivalent$ (row): $P = \Identity{U}$, i.e.\ the subset $\Set{\Identity{U}} \times \GL{V}$;
        \item $\DomainEquivalent$ (column): $Q = \Identity{V}$, i.e.\ the subset $\GL{U} \times \Set{\Identity{V}}$;
        \item $\Similar$ (when $U = V$): $P = Q$, i.e.\ the diagonal $\Set{(R, R) \mid R \in \GL{V}}$, giving
              the conjugation $T = \Composition{Q}{\Composition{S}{Q^{-1}}}$.
    \end{itemize}
    Each subset contains $(\Identity{U}, \Identity{V})$ and is closed under inverses and composition, so the
    three displayed witnesses stay inside it. The same computation therefore proves each of the four
    relations reflexive, symmetric, and transitive.
\end{proof}

Since the matrix relations are nothing but these map relations restricted to coordinate spaces, the
content still to be settled is how they connect a \emph{fixed abstract} map to its coordinate matrices in
\emph{arbitrary} bases. The next two theorems are that bridge: a map relation holds iff the corresponding
relation holds between coordinate matrices, in \emph{any} bases; and conversely two matrices are related
iff they present a single map.

\begin{theorem}[Equivalence in coordinates]
\label{thm:matrices-of-map-equivalent}
    Let $S, T : U \to V$ be linear maps.
    \begin{enumerate}
        \item For \emph{any} ordered bases $\mathcal B_U$ of $U$ and $\mathcal B_V$ of $V$,
              \[
                  S \Equivalent T
                  \quad\Longleftrightarrow\quad
                  \CoordMatrix[\mathcal B_U][\mathcal B_V]{S} \Equivalent \CoordMatrix[\mathcal B_U][\mathcal B_V]{T} .
              \]
        \item Two matrices $A, B \in \CFMat{Y}{X}[K]$ are equivalent iff they are the matrices of one
              common map: $A = \CoordMatrix[\mathcal B_U][\mathcal B_V]{R}$ and
              $B = \CoordMatrix[\mathcal C_U][\mathcal C_V]{R}$ for some $R : U \to V$ and ordered bases.
    \end{enumerate}
\end{theorem}

\begin{proof}
    \emph{(1).} If $T = \Composition{Q}{\Composition{S}{P^{-1}}}$ with $P \in \GL{U}$, $Q \in \GL{V}$,
    then the composition theorem (Theorem~\ref{thm:composition-matrix}) and the inversion rule
    (Theorem~\ref{thm:inverse-matrix}) give
    \[
        \CoordMatrix[\mathcal B_U][\mathcal B_V]{T}
        = \Composition{\CoordMatrix[\mathcal B_V][\mathcal B_V]{Q}}
            {\Composition{\CoordMatrix[\mathcal B_U][\mathcal B_V]{S}}
              {\bigl(\CoordMatrix[\mathcal B_U][\mathcal B_U]{P}\bigr)^{-1}}} ,
    \]
    an equivalence of the two matrices via the invertible factors
    $\CoordMatrix[\mathcal B_V][\mathcal B_V]{Q}$ and $\CoordMatrix[\mathcal B_U][\mathcal B_U]{P}$
    (Corollary~\ref{cor:invertible-detect}). Conversely, if
    $\CoordMatrix[\mathcal B_U][\mathcal B_V]{S}$ and $\CoordMatrix[\mathcal B_U][\mathcal B_V]{T}$ are
    equivalent matrices, their invertible factors are --- again by
    Corollary~\ref{cor:invertible-detect} --- the matrices
    $\CoordMatrix[\mathcal B_V][\mathcal B_V]{Q}$ and $\CoordMatrix[\mathcal B_U][\mathcal B_U]{P}$ of
    automorphisms $Q \in \GL{V}$ and $P \in \GL{U}$; reading the display backwards gives
    $T = \Composition{Q}{\Composition{S}{P^{-1}}}$, i.e.\ $S \Equivalent T$.

    \emph{(2).} ($\Leftarrow$) The change-of-bases formula (Corollary~\ref{cor:change-of-bases}) gives
    $A = \CoordMatrix[\mathcal C_V][\mathcal B_V]{\Identity{V}}\, B\, \CoordMatrix[\mathcal B_U][\mathcal C_U]{\Identity{U}}$,
    a product $Q B P^{-1}$ of $B$ with invertible transition matrices, so $A \Equivalent B$.
    ($\Rightarrow$) If $A = Q B P^{-1}$, take $R = B : K^{(X)} \to K^{(Y)}$ with $\mathcal C_U, \mathcal C_V$
    the canonical bases (so $\CoordMatrix[\mathcal C_U][\mathcal C_V]{R} = B$); the columns of $P^{-1}$ and
    $Q^{-1}$ are ordered bases $\mathcal B_U, \mathcal B_V$ with
    $\CoordMatrix[\mathcal B_U][\mathcal C_U]{\Identity{}} = P^{-1}$ and
    $\CoordMatrix[\mathcal C_V][\mathcal B_V]{\Identity{}} = Q$, whence
    $\CoordMatrix[\mathcal B_U][\mathcal B_V]{R} = Q B P^{-1} = A$.
\end{proof}

\begin{theorem}[Similarity in coordinates]
\label{thm:matrices-of-endo-similar}
    Let $S, T \in \End[K]{V}$ be endomorphisms.
    \begin{enumerate}
        \item For \emph{any} ordered basis $\mathcal B$ of $V$,
              \[
                  S \Similar T
                  \quad\Longleftrightarrow\quad
                  \CoordMatrix[\mathcal B][\mathcal B]{S} \Similar \CoordMatrix[\mathcal B][\mathcal B]{T} .
              \]
        \item Square matrices $A, B \in \CFMat{X}{X}[K]$ are similar iff they are the matrices of one
              common endomorphism: $A = \CoordMatrix[\mathcal B][\mathcal B]{R}$ and
              $B = \CoordMatrix[\mathcal C][\mathcal C]{R}$ for some $R \in \End[K]{V}$ and ordered bases
              $\mathcal B, \mathcal C$ of $V$.
    \end{enumerate}
\end{theorem}

\begin{proof}
    Both parts repeat the previous proof with one basis on each side and $P = Q$. For (1), if
    $T = \Composition{Q^{-1}}{\Composition{S}{Q}}$ then
    $\CoordMatrix[\mathcal B][\mathcal B]{T} = \bigl(\CoordMatrix[\mathcal B][\mathcal B]{Q}\bigr)^{-1}\, \CoordMatrix[\mathcal B][\mathcal B]{S}\, \CoordMatrix[\mathcal B][\mathcal B]{Q}$
    with $\CoordMatrix[\mathcal B][\mathcal B]{Q}$ invertible, and conversely. For (2), the
    ($\Leftarrow$) direction is Corollary~\ref{cor:change-of-bases} with a single basis per side,
    yielding $A = \CoordMatrix[\mathcal C][\mathcal B]{\Identity{V}}\, B\, \CoordMatrix[\mathcal B][\mathcal C]{\Identity{V}} = Q^{-1} B Q$;
    the ($\Rightarrow$) direction takes $R = B$ on $K^{(X)}$ with $\mathcal C$ canonical and $\mathcal B$
    the columns of $Q$.
\end{proof}

The sharp statement is at the level of the \emph{subspaces} themselves, and it splits cleanly along the
two one-sided relations: a row (codomain) change leaves the kernel untouched, a column (domain) change
leaves the range untouched. The numerical invariants --- rank and nullity --- are then mere shadows.

\begin{proposition}[Kernel is codomain-invariant; range is domain-invariant]
\label{prop:ker-rng-one-sided-invariant}
    Let $A, B \in \CFMat{Y}{X}[K]$.
    \begin{enumerate}
        \item If $A \CodomainEquivalent B$ then $\Ker{A} = \Ker{B}$: the kernel depends only on the
              codomain (row) class.
        \item If $A \DomainEquivalent B$ then $\Rng{A} = \Rng{B}$: the range depends only on the domain
              (column) class.
    \end{enumerate}
\end{proposition}

\begin{proof}
    \emph{(1).} Here $A = Q B$ with $Q$ invertible, in particular injective; so for every
    $x \in K^{(X)}$, $A(x) = Q(B(x)) = 0 \iff B(x) = 0$, that is $\Ker{A} = \Ker{B}$.

    \emph{(2).} Here $A = B P^{-1}$ with $P^{-1}$ invertible, in particular surjective; so
    \[
        \Rng{A} = \DirectImage{A}{K^{(X)}}
        = \DirectImage{B}{\DirectImage{P^{-1}}{K^{(X)}}}
        = \DirectImage{B}{K^{(X)}} = \Rng{B} . \qedhere
    \]
\end{proof}

\begin{corollary}[Rank and nullity are equivalence invariants]
\label{prop:rank-equiv-invariant}
    If $A \Equivalent B$ then $\Rank{A} = \Rank{B}$ and $\Nullity{A} = \Nullity{B}$.
\end{corollary}

\begin{proof}
    By Proposition~\ref{prop:equivalence-factors} factor the equivalence as a column step followed by a
    row step: there is $C$ with $B \DomainEquivalent C$ and $C \CodomainEquivalent A$, namely
    $C = B P^{-1}$ and $A = Q C$.
    \begin{itemize}
        \item \emph{Rank.} The column step fixes the range exactly,
              $\Rng{C} = \Rng{B}$ (Proposition~\ref{prop:ker-rng-one-sided-invariant}(2)); the row step
              $A = Q C$ moves it by the isomorphism $Q$, $\Rng{A} = \DirectImage{Q}{\Rng{C}}$, preserving
              its dimension. Hence $\Rank{A} = \Rank{C} = \Rank{B}$.
        \item \emph{Nullity.} The row step fixes the kernel exactly,
              $\Ker{A} = \Ker{C}$ (Proposition~\ref{prop:ker-rng-one-sided-invariant}(1)); the column
              step $C = B P^{-1}$ moves it by the isomorphism $P$, $\Ker{C} = \DirectImage{P}{\Ker{B}}$,
              preserving its dimension. Hence $\Nullity{A} = \Nullity{C} = \Nullity{B}$.
    \end{itemize}
\end{proof}

\begin{remark}[The basis-free content: the equivalence and similarity classes]
\label{rem:basis-free-content}
    Theorem~\ref{thm:matrices-of-map-equivalent} says the set of all matrices of a fixed map $T : U \to V$
    is exactly one equivalence class; for an endomorphism, insisting on one basis for both source and
    target (Theorem~\ref{thm:matrices-of-endo-similar}) refines this to a single similarity class. So a
    quantity computed from a matrix is an intrinsic invariant of the underlying map precisely when it
    depends only on the class, not on the chosen representative. We have already isolated one invariant
    of the equivalence class --- the rank (Corollary~\ref{prop:rank-equiv-invariant}), which by the
    forthcoming normal-form theorem turns out to be a \emph{complete} invariant of equivalence. The
    similarity classes are finer, and their classification needs invariants special to the endomorphism
    case --- the \textbf{determinant}, the \textbf{trace}, and the \textbf{characteristic polynomial}.
    These are constructed in later chapters and there shown to depend only on the similarity class; we
    name them here only to signpost where the theory is heading.
\end{remark}

\subsection{Worked Example}

\begin{example}[Two matrices of one endomorphism are similar]
\label{ex:similar-numeric}
    Return to Example~\ref{ex:change-of-bases-numeric}: $T(x, y) = (2x + y,\; x)$ on $\bb{R}^2$, with
    $\mathcal B$ standard and $\mathcal C = (\CanonicalVector{0} + \CanonicalVector{1},\, \CanonicalVector{1})$.
    There we found
    $\CoordMatrix[\mathcal B][\mathcal B]{T} = \begin{psmallmatrix} 2 & 1 \\ 1 & 0 \end{psmallmatrix}$ and
    $\CoordMatrix[\mathcal C][\mathcal C]{T} = \begin{psmallmatrix} 3 & 1 \\ -2 & -1 \end{psmallmatrix}$.
    With $Q = \CoordMatrix[\mathcal B][\mathcal C]{\Identity{}} = \begin{psmallmatrix} 1 & 0 \\ -1 & 1 \end{psmallmatrix}$
    and $Q^{-1} = \begin{psmallmatrix} 1 & 0 \\ 1 & 1 \end{psmallmatrix}$,
    \[
        Q^{-1}\, \CoordMatrix[\mathcal C][\mathcal C]{T}\, Q
        = \begin{pmatrix} 1 & 0 \\ 1 & 1 \end{pmatrix}
          \begin{pmatrix} 3 & 1 \\ -2 & -1 \end{pmatrix}
          \begin{pmatrix} 1 & 0 \\ -1 & 1 \end{pmatrix}
        = \begin{pmatrix} 2 & 1 \\ 1 & 0 \end{pmatrix}
        = \CoordMatrix[\mathcal B][\mathcal B]{T} ,
    \]
    so $\CoordMatrix[\mathcal B][\mathcal B]{T} \Similar \CoordMatrix[\mathcal C][\mathcal C]{T}$ ---
    exactly part~(2) of Theorem~\ref{thm:matrices-of-endo-similar}, the two matrices of the one
    endomorphism $T$. Both have rank $2$
    (Corollary~\ref{prop:rank-equiv-invariant}); looking ahead, they will also share the finer
    similarity invariants --- determinant, trace, characteristic polynomial --- once those are defined.
\end{example}
